\section{Lower mass and solvability for all measures}
\label{sec:lower-mass-solvability}

Throughout the rest of the chapter, $h$ denotes a Hermitian form on $X$, used as the background; the sections that use another background say so explicitly. No singularity type is prescribed: the unknown potentials range over all of $\PSH(X,h)$. Constants belong to $\PSH(X,h)$, and every $u\in\PSH(X,h)$ is bounded above, so $u\preceq v$ for every bounded $v\in\PSH(X,h)$.

The masses of bounded potentials are measured by the lower and upper masses $v_\pm(h)$ of \cref{def:bmp}, the infimum and the supremum of $\int_X(h+\ddc v)^n$ over the bounded $v\in\PSH(X,h)$; recall that $h$ has positive lower mass if $v_-(h)>0$. If $h$ is K\"ahler, Stokes' theorem gives $v_-(h)=v_+(h)=\int_Xh^n$; on a non-K\"ahler manifold the lower mass can vanish (\cref{thm:vanishing-lower-mass}).

For a nonzero nonnegative $q\in L^p(\mathrm{d}V)$ with $p>1$, \cref{thm:kn-continuous} gives a constant $c>0$ and a continuous $u\in\PSH(X,h)$ with $\mathrm{MA}_h(u)=cq\,\mathrm{d}V$, and \cref{prop:vanishing-exponent} shows that $c$ is determined by $q$. We call $c_h(q)=c$ the \emph{normalization scalar}\index{Normalization scalar!definition} of $q$, and for the measure $\nu=q\,\mathrm{d}V$ we also write $c_h(\nu)=c_h(q)$. We say that a finite measure $\nu$ admits a \emph{positive scaled subsolution} if there are $u\in\PSH(X,h)$ and a constant $a>0$ such that $\mathrm{MA}^{\mathrm{ext}}_h(u)\geq a\nu$.

If $\nu$ admits a solution of $\mathrm{MA}^{\mathrm{ext}}_h(u)=c\nu$ with $c>0$, then it admits in particular a positive scaled subsolution. We show that, when $h$ has zero lower mass, some probability density admits no positive scaled subsolution. Conversely, if $v_-(h)>0$, every nonzero finite non-pluripolar measure admits such a solution (\cref{cor:positive-lower-mass-solvability} in \cref{chap:relative-finite-measures}). The key estimate is a consequence of strict domination: if $\mathrm{MA}^{\mathrm{ext}}_h(u)\geq aq\,\mathrm{d}V$, then $a\leq c_h(q)$. When $v_-(h)=0$, there are smooth probability densities with arbitrarily small normalization scalars, and a weighted sum of such densities, with scalars decreasing fast enough, is a single probability density admitting no positive scaled subsolution.

\begin{lemma}\label{lem:bounded-normalizer-subsolution}\label{lem:smooth-normalizer-subsolution-bound}
Let $0\leq q\in L^p(\mathrm{d}V)$, $p>1$, be nonzero. If $u\in\PSH(X,h)$ and $a>0$ satisfy
\[
 \mathrm{MA}^{\mathrm{ext}}_h(u)\geq aq\,\mathrm{d}V,
\]
then $a\leq c_h(q)$. Consequently, for nonzero nonnegative $q,r\in L^p(\mathrm{d}V)$,
\[
 q\leq r\ \Longrightarrow\ c_h(q)\geq c_h(r),
 \qquad c_h(tq)=t^{-1}c_h(q)\quad(t>0).
\]
\end{lemma}
\begin{proof}
Let $v$ be a continuous solution of $\mathrm{MA}_h(v)=c_h(q)q\,\mathrm{d}V$, and suppose, to reach a contradiction, that $a>c_h(q)$. Then $\tau=c_h(q)/a<1$ and, for every real $A$,
\[
 \mathrm{MA}_h(v+A)=c_h(q)q\,\mathrm{d}V
       \leq\tau\,\mathrm{MA}^{\mathrm{ext}}_h(u).
\]
The translate $v+A$ is bounded, so $P_{h/2}(v+A)\geq\inf_Xv+A$ is not identically $-\infty$, and $u\preceq v+A$. \Cref{thm:singular-strict-domination}, applied with $\delta=1/2$ and $c=\tau$ to $v+A$ in place of $u$ and to $u$ in place of $v$, therefore gives $v+A\geq u$ for every real $A$. Letting $A\to-\infty$ at a point where $u$ is finite gives a contradiction. Hence $a\leq c_h(q)$.

If $q\leq r$, a continuous solution of $\mathrm{MA}_h(w)=c_h(r)r\,\mathrm{d}V$ satisfies $\mathrm{MA}_h(w)\geq c_h(r)q\,\mathrm{d}V$, so the first assertion gives $c_h(r)\leq c_h(q)$. Finally, the equation $\mathrm{MA}_h(v)=c_h(q)q\,\mathrm{d}V$ can be written as $\mathrm{MA}_h(v)=t^{-1}c_h(q)\,(tq)\,\mathrm{d}V$, and uniqueness of the scalar gives $c_h(tq)=t^{-1}c_h(q)$.
\end{proof}

\begin{lemma}\label{lem:lp-normalizer-continuity}
Fix $p>1$. If nonzero nonnegative densities $q_j,q\in L^p(\mathrm{d}V)$ satisfy $q_j\to q$ in $L^p$, then $c_h(q_j)\to c_h(q)$.
\end{lemma}
\begin{proof}
We first show that the scalars $c_h(q_j)$ lie in a compact subset of $(0,+\infty)$. Fix a Gauduchon metric $\gamma$ on $X$ (\cref{def:gauduchon-metric}; in dimension one, any Hermitian metric). Pairing the equation for $q_j$ with $\gamma^{n-1}$, as in \cref{prop:big-singular-gauduchon-trace} with $\theta=h$ and barrier $\rho=0$, gives
\[
 c_h(q_j)^{1/n}\int_X\left(\frac{q_j\,\mathrm{d}V}{\gamma^n}\right)^{1/n}\gamma^n
 \leq\int_Xh\wedge\gamma^{n-1}.
\]
Since $q_j^{1/n}\to q^{1/n}$ in $L^n$, the integrals on the left converge to the corresponding integral for $q$, which is positive. Hence the scalars are bounded above. The lower-density estimate of \cref{thm:lp-potential-estimate} gives $c_h(q_j)\|q_j\|_{L^p}\geq c_p>0$, so they are also bounded away from zero.

Consequently, the normalized continuous solutions $u_j$ of $\mathrm{MA}_h(u_j)=c_h(q_j)q_j\,\mathrm{d}V$ have uniformly bounded $L^p$ densities. By \cref{thm:lp-uniform-compactness}, every subsequence has a further subsequence along which $u_j$ converges uniformly to some $u$ and $c_h(q_j)$ converges to some $c>0$. Continuity of bounded Monge--Amp\`ere products under uniform convergence gives $\mathrm{MA}_h(u)=cq\,\mathrm{d}V$, and uniqueness of the scalar gives $c=c_h(q)$. Since every subsequence of $(c_h(q_j))$ has a further subsequence converging to $c_h(q)$, the whole sequence converges to $c_h(q)$.
\end{proof}

\begin{theorem}\label{thm:lower-mass-obstructing-measure}\index{Non-pluripolar measure!solvability obstruction}
Suppose $v_-(h)=0$. There is a probability measure $\mu=f\,\mathrm{d}V$, with $f\in L^1(\mathrm{d}V)$ and $f\geq m$ almost everywhere for a constant $m>0$, such that no $u\in\PSH(X,h)$ and $a>0$ satisfy
\begin{equation}\label{eq:lower-mass-impossible-subsolution}
 \mathrm{MA}^{\mathrm{ext}}_h(u)\geq a\mu.
\end{equation}
In particular, the equation $\mathrm{MA}^{\mathrm{ext}}_h(u)=c\mu$ has no solution $u\in\PSH(X,h)$ for any $c>0$; neither a singularity type nor finiteness of the mass of $u$ is imposed.
\end{theorem}
\begin{proof}
By \cref{thm:gaud-volume-mass-extrema,thm:gaud-volume-dual-formula}, $v_-(h)$ is the infimum of the normalization scalars $c_h(\nu)$ over smooth strictly positive probability measures $\nu$. Since $v_-(h)=0$, we can choose such measures $\nu_j=f_j\,\mathrm{d}V$ with
\[
 c_h(\nu_j)\leq 2^{-2j}\qquad(j\geq1).
\]
Set
\[
 \mu=\sum_{j=1}^{\infty}2^{-j}\nu_j,
 \qquad f=\sum_{j=1}^{\infty}2^{-j}f_j.
\]
By monotone convergence, $\int_Xf\,\mathrm{d}V=1$, so $f\in L^1$. Moreover $f\geq\tfrac12\min_X f_1=:m>0$, and $\mu$ is non-pluripolar, being absolutely continuous with respect to $\mathrm{d}V$.

Suppose that \eqref{eq:lower-mass-impossible-subsolution} held for some $u$ and $a>0$. Since $\mu\geq2^{-j}\nu_j$, we would have
\[
 \mathrm{MA}^{\mathrm{ext}}_h(u)\geq a2^{-j}\nu_j
\]
for every $j$, and \cref{lem:bounded-normalizer-subsolution} would give $a2^{-j}\leq c_h(\nu_j)\leq2^{-2j}$. Thus $a\leq2^{-j}$ for every $j$, contradicting $a>0$.
\end{proof}

\Needspace{12\baselineskip}
\begin{corollary}\label{cor:lower-mass-lone-instability}
Assume $v_-(h)=0$, and let $\mu=f\,\mathrm{d}V$ be the measure of \cref{thm:lower-mass-obstructing-measure}. Then $f\notin L^p(\mathrm{d}V)$ for every $p>1$. Moreover, for every probability density $g\in L^1(\mathrm{d}V)$ and every $0<t<1$, the probability density
\[
 g_t=(1-t)g+tf
\]
admits no positive scaled subsolution, although $\|g_t-g\|_{L^1}\leq2t$. In particular, the probability densities admitting no positive scaled subsolution are dense among all probability densities in $L^1(\mathrm{d}V)$.
\end{corollary}
\begin{proof}
If $f\in L^p$ for some $p>1$, \cref{thm:kn-continuous} would give a continuous solution of $\mathrm{MA}_h(u)=cf\,\mathrm{d}V$ with $c>0$, contradicting the theorem. Since $g_t\,\mathrm{d}V\geq t\mu$, a positive scaled subsolution for $g_t\,\mathrm{d}V$ would also give one for $\mu$. Finally,
$\|g_t-g\|_{L^1}=t\|f-g\|_{L^1}\leq2t$.
\end{proof}

Smooth strictly positive probability densities are also dense in $L^1(\mathrm{d}V)$, and each of them admits a smooth solution of the normalized equation. Thus, when the lower mass vanishes, solvability is not preserved under $L^1$ approximation. The obstruction comes from the way the density concentrates, not from its size: the density in \cref{thm:lower-mass-obstructing-measure} is bounded below by a positive constant. It occurs, for instance, on the compact threefold of \cref{thm:vanishing-lower-mass}, where every Hermitian form has zero lower mass.

\Cref{thm:lower-mass-obstructing-measure} does not exclude unbounded solutions of the exponential equation $\mathrm{MA}^{\mathrm{ext}}_h(u)=\mathrm{e}^{\lambda u}\mu$ with $\lambda>0$, because the factor $\mathrm{e}^{\lambda u}$ need not have a positive lower bound. It does exclude bounded solutions, since a bounded solution would satisfy $\mathrm{MA}^{\mathrm{ext}}_h(u)\geq\mathrm{e}^{\lambda\inf_Xu}\mu$. On suitable products, \cref{thm:zero-threshold-solvable-concentration} chooses the concentrating components of such a density so that its exponential equation has a solution for every positive exponent. We do not know whether this holds for every density admitting no positive scaled subsolution, nor which individual measures admit a solution of the normalized equation when $v_-(h)=0$.
